% This template is dedicated to the public domain under CC0 1.0: use it for
% anything, no credit needed. https://creativecommons.org/publicdomain/zero/1.0/

\documentclass[11pt]{article}
\usepackage[margin=1in]{geometry}
\usepackage{amsmath}
\usepackage{booktabs}
\usepackage{tikz}
\usepackage[backend=bibtex]{biblatex}
\addbibresource{references.bib}
\usepackage{hyperref}
\usepackage{maleficium-doc}   % the Maleficium look; remove for plain LaTeX
\usepackage{maleficium-footer}   % "Made with Maleficium" mark; delete this line to remove it

\title{The Title of Your Article}
\author{Your Name \\ \small Your Institution}
\date{\today}

\begin{document}
\maketitle

\begin{abstract}
One paragraph that says what you did, why it matters and what you found.
\end{abstract}

\section{Introduction}
\label{sec:intro}
State the problem and your contribution. Cite prior work like
this~\cite{knuth1984}, and point ahead to Section~\ref{sec:method}.

\section{Method}
\label{sec:method}
Number the equations you refer to:
\begin{equation}
  \label{eq:energy}
  E = m c^2 .
\end{equation}
Equation~\eqref{eq:energy} is then easy to point back to.

\section{Results}
Table~\ref{tab:results} and Figure~\ref{fig:sketch} summarise the outcome.

\begin{table}[h]
  \centering
  \begin{tabular}{lrr}
    \toprule
    Method   & Accuracy & Time (s) \\
    \midrule
    Baseline & 0.81     & 12.0     \\
    Ours     & 0.93     & 9.5      \\
    \bottomrule
  \end{tabular}
  \caption{Replace these numbers with yours.}
  \label{tab:results}
\end{table}

\begin{figure}[h]
  \centering
  \begin{tikzpicture}
    \draw[->] (0,0) -- (4,0) node[right] {$x$};
    \draw[->] (0,0) -- (0,2.5) node[above] {$y$};
    \draw[thick, mlaccent] (0,0.2) .. controls (1.5,2.4) and (2.5,0.4) .. (4,2);
  \end{tikzpicture}
  \caption{A figure drawn in TikZ; use \texttt{\textbackslash includegraphics} for images.}
  \label{fig:sketch}
\end{figure}

\section{Conclusion}
Say what the reader should remember.

\printbibliography
\end{document}
